%杨舒云的实验报告编辑界面，使用了Huanyu Shi，2019级的模板，杨舒云在此拜谢ORZ!

%!TEX program = xelatex
\documentclass[dvipsnames, svgnames,a4paper,11pt]{article}
\input{YsySettings} 
\usepackage{lipsum}
\usepackage{adjustbox}
\usepackage{diagbox}

\begin{document}
	
	
	
	%实验报告封面	
	\begin{table}
		\renewcommand\arraystretch{1.7}
		\begin{tabularx}{\textwidth}{
				|X|X|X|X
				|X|X|X|X|}
			\hline
			\multicolumn{2}{|c|}{预习报告}&\multicolumn{2}{|c|}{实验记录}&\multicolumn{2}{|c|}{分析讨论}&\multicolumn{2}{|c|}{总成绩}\\
			\hline
			\LARGE20 & & \LARGE30 & & \LARGE30 & & \LARGE80 & \\
			\hline
		\end{tabularx}
	\end{table}
	
	\begin{table}
		\renewcommand\arraystretch{1.7}
		\begin{tabularx}{\textwidth}{|X|X|X|X|}
			\hline
			年级、专业： & 2022级 物理学 &组号： & 2\\
			\hline
			姓名： & 杨舒云  & 学号： & 22344020\\
			\hline
			实验时间： & 2023/12/14 & 教师签名： & \\
			\hline
		\end{tabularx}
	\end{table}
	
	\begin{center}
		\LARGE BA3 \quad 应用AMR测量磁场分布
	\end{center}
	
	\textbf{【实验报告注意事项】}
	\begin{enumerate}
		\item 实验报告由三部分组成：
		\begin{enumerate}
			\item 预习报告：课前认真研读实验讲义，弄清实验原理；实验所需的仪器设备、用具及其使用、完成课前预习思考题；了解实验需要测量的物理量，并根据要求提前准备实验记录表格（可以参考实验报告模板，可以打印）。\textcolor{red}{\textbf{（25分）}}
			\item 实验记录：认真、客观记录实验条件、实验过程中的现象以及数据。实验记录请用珠笔或者钢笔书写并签名（\textcolor{red}{\textbf{用铅笔记录的被认为无效}}）。\textcolor{red}{\textbf{保持原始记录，包括写错删除部分，如因误记需要修改记录，必须按规范修改。}}（不得输入电脑打印，但可扫描手记后打印扫描件）；离开前请实验教师检查记录并签名。\textcolor{red}{\textbf{（30分）}}
			\item 数据处理及分析讨论：处理实验原始数据（学习仪器使用类型的实验除外），对数据的可靠性和合理性进行分析；按规范呈现数据和结果（图、表），包括数据、图表按顺序编号及其引用；分析物理现象（含回答实验思考题，写出问题思考过程，必要时按规范引用数据）；最后得出结论。\textcolor{red}{\textbf{（25分）}}
		\end{enumerate}
		\textbf{实验报告就是将预习报告、实验记录、和数据处理与分析合起来，加上本页封面。\textcolor{red}{（80分）}}
		\item 每次完成实验后的一周内交\textbf{实验报告}（特殊情况不能超过两周）。
	\end{enumerate}
	
	\textbf{【实验安全注意事项】}	
	\begin{enumerate}
		\item 禁止将实验仪处于强磁场中，否则会严重影响实验结果;
		\item 为了降低实验仪间磁场的相互干扰，任意两台实验仪之间的距离应大于 3 米;
		\item 实验前请先调水平实验仪; 
		\item 在操作所有的手动调节螺钉时应用力适度，以免滑丝; 
		\item 为保证使用安全，三芯电源须可靠接地。		
	\end{enumerate}
	
	\textbf{【特别鸣谢及模板说明】}	
	
	感谢2019级学长石寰宇为本实验报告提供\LaTeX 模板。
	
	
	
	\clearpage
	\tableofcontents
	\clearpage
	
	
	
	%预习报告	
	\setcounter{section}{0}
	\section{BA3 应用AMR测量磁场分布 \quad\heiti 预习报告}
	
	\subsection{实验目的}
	\begin{enumerate}
		\item 学习各向异性磁阻与磁传感器原理；
		\item 学习磁场的测量方法；
		\item 测量磁阻传感器的磁电转换特性和各向异性特性；
		\item 测量地磁场强度、磁倾角、磁偏角。
	\end{enumerate}
	
	\subsection{仪器用具}
	\begin{table}[htbp]
		\centering
		\renewcommand\arraystretch{1.6}
		% \setlength{\tabcolsep}{10mm}
		\begin{tabular}{p{0.05\textwidth}|p{0.20\textwidth}|p{0.05\textwidth}|p{0.5\textwidth}}
			\hline
			编号& 仪器用具名称 & 数量 &  主要参数（型号，测量范围，测量精度等） \\
			\hline
			1& 各项异性磁阻传感器 & 1台 & 含亥姆霍兹线圈、电源、AMR测量装置 \\
			
			2& 电流输出线 & 3根 & -- \\
			
			3& 水准仪 & 1 个 & -- \\
			\hline
		\end{tabular}
	\end{table}
	
	\begin{figure}[htbp]
		\centering
		\includegraphics[width=0.7\textwidth]{BA3GraADD1.jpg}
		\caption{实验仪器}
		\label{fig:figADD1}
	\end{figure}
	
	\textbf{磁阻传感器的输出信号通常须经放大电路放大后，再接显示电路，故由显示电压计算磁 场强度时还需考虑放大器的放大倍数。本实验仪电桥工作电压5V，放大器放大倍数50，磁感应强度为1Gauss时，对应的输出电压为0.25V。}
	
	\subsection{原理概述}
	引言：物质在磁场中电阻率发生变化的现象称为磁阻效应，利用这一效应，可通过测量材料的电阻(及其变化)反过来测量磁场。磁阻效应材料发展很快，磁阻元件的发展分别经历了磁阻(MR)，各向异性磁阻(AMR)，巨磁阻(GMR)，庞磁阻(CMR)等阶段。
	磁阻传感器可用于直接测量磁场或磁场变化，如弱磁场测量，地磁场测量，各种导航系统中的罗盘，计算机中的磁盘驱动器，各种磁卡机等等。也可通过磁场变化测量其它物理量，如利用磁阻效应已制成各种位移、角度、转速传感器，各种接近开关，隔离开关，广泛用于汽车，家电及各类需要自动检测与控制的领域。
	本实验利用 AMR 传感器测量地磁场及亥姆霍兹线圈磁场分布，并探究磁介质对其周围空间磁场分布的影响。对于学过《电磁学》的同学，这部分的实验内容是对理论的验证。
	
	\subsubsection{简述各向异性磁阻特性与磁传感器原理（4分）}
	各向异性磁阻（AMR）特性是指材料的电阻值随着外部磁场的方向和强度的变化而改变。这种特性是由材料内部电子的自旋与磁场的相互作用导致的。AMR效应在许多磁性材料中都存在，尤其在铁磁性材料中表现得更为明显。
	
	\begin{figure}[htbp]
		\centering
		\includegraphics[width=0.2\textwidth]{BA3GraAD1.jpg}
		\caption{各向异性磁阻（参考文献[1]）}
		\label{fig:figAD1}
	\end{figure}
	
	各向异性磁阻特性：
	\begin{enumerate}
		\item \textbf{依赖于磁场方向：} 当磁场方向改变时，材料的电阻值也会发生变化。这是因为电子在材料中的运动受到磁场方向的影响。
		
		\item \textbf{磁场强度的影响：} 随着磁场强度的增加，电阻值的变化也会增大，直到达到饱和点。
		
		\item \textbf{温度效应：} AMR效应也受温度的影响，不同温度下材料的电阻率和磁阻变化特性可能不同。
	\end{enumerate}
	
	磁传感器原理：利用AMR效应的磁传感器主要是测量材料电阻的变化来检测磁场的存在、强度和方向。
	\begin{enumerate}
		\item \textbf{检测原理：} 当磁场存在时，磁传感器内的AMR材料会因为自旋与磁场的相互作用而改变电阻值。通过精确测量这种电阻的变化，可以推断出磁场的性质。
		
		\item \textbf{输出信号：} 磁传感器通常将电阻变化转换为电压或电流的变化，以便于读取和处理。
		
		\item \textbf{应用领域：} AMR磁传感器广泛应用于位置和速度传感、磁场强度测量、磁盘驱动器的数据读取等领域。
	\end{enumerate}
	
	AMR磁传感器的关键优势在于它们能提供关于磁场方向的详细信息，并且相比其他类型的磁传感器（如霍尔效应传感器）通常具有更高的精度和灵敏度。然而，它们也可能受到温度变化和其他环境因素的影响，因此在设计和使用时需要考虑这些因素。
	
	
	\subsubsection{简述亥姆霍兹线圈产生磁场的特点（4分）}
	如\cref{fig:figAD2},亥姆霍兹线圈是由两个相同的、平行的、同轴的圆形线圈组成的装置，当两个线圈都有相同电流流过时，它们在其中心区域产生一个近似均匀的磁场。这种线圈的设计是为了在一个相对较大的空间区域内提供一个均匀的磁场，这在物理学和工程学实验中是非常有用的。
	
	\begin{figure}[htbp]
		\centering
		\includegraphics[width=0.2\textwidth]{BA3GraAD2.jpg}
		\caption{亥姆霍兹线圈（参考文献[1]）}
		\label{fig:figAD2}
	\end{figure}
	
	亥姆霍兹线圈磁场的特点：
	\begin{enumerate}
		\item \textbf{均匀性：} 在两个线圈的中点附近，磁场非常均匀。这是亥姆霍兹线圈最显著的特点，使其成为实验中理想的磁场源。
		
		\item \textbf{可调性：} 通过改变流经线圈的电流，可以调节磁场的强度。这使得亥姆霍兹线圈适用于需要不同磁场强度的实验。
		
		\item \textbf{方向性：} 亥姆霍兹线圈产生的磁场方向与线圈平面垂直，这对于需要特定磁场方向的应用非常有用。
		
		\item \textbf{设计简单：} 亥姆霍兹线圈结构简单，制作相对容易，能够满足多种实验需求。
		
		\item \textbf{较大的有效区域：} 相比单个线圈，亥姆霍兹线圈在其中心附近提供了一个更大的区域，其中磁场保持均匀。
		
		\item \textbf{应用广泛：} 用于物理实验（如电子的荷质比实验）、生物医学（如磁共振成像MRI的研究）、精密测量磁场等。
	\end{enumerate}
	
	这些特点使亥姆霍兹线圈成为理想的实验工具，特别是在需要精确控制磁场均匀性和强度的情况下。
	
	亥姆霍兹线圈的数学描述主要涉及计算线圈产生的磁场。亥姆霍兹线圈由两个同半径 \( R \) 的圆形线圈组成，它们沿同一轴线排列，彼此间距也是 \( R \)。当两个线圈中流过相同的电流 \( I \) 时，它们在中心点产生近似均匀的磁场。
	
	在线圈中心点（\( x = 0 \)）的磁场强度 \( B \) 可以用以下公式计算：
	
	\begin{equation*}
		B = \left( \frac{4}{5} \right)^{3/2} \frac{\mu_0 N I}{R}
	\end{equation*}
	
	其中，$\mu_0 \text{ 是真空磁导率，约等于 } 4\pi \times 10^{-7} \, \text{T}\cdot\text{m/A},  N \text{ 是每个线圈的匝数},  I \text{ 是通过线圈的电流}$。 采用国际单位制时，由上式计算出的磁感应强度单位为特斯拉($1 T=10000 Gauss$)。本实验仪N=310，R=0.14m，线圈电流为1mA时，赫姆霍兹线圈中部的磁感应强度为0.02Gauss。
	
	这个方程表明，中心点的磁场强度与线圈的半径、匝数以及流经线圈的电流成正比。亥姆霍兹线圈设计的核心在于使线圈间距与半径相等，这样可以在中间区域产生一个尽可能均匀的磁场。
	
	
	\subsubsection{简述地磁场特性（4分）}
	地磁场是地球周围的自然磁场，具有以下主要特性：
	
	\begin{enumerate}
		\item \textbf{源头：} 地磁场主要由地球内核的熔融金属流动产生，尤其是由于外核的液态铁的对流运动。
		
		\item \textbf{形态：} 在大致上，地磁场类似于一个倾斜的偶极磁体，其磁极大约与地球的地理极相近。
		
		\item \textbf{强度：} 地磁场的强度不均匀，从赤道到两极逐渐增强。平均强度约为25-65微特斯拉（μT）。
		
		\item \textbf{方向：} 地磁场的方向（磁偏角和磁倾角）随地理位置变化。磁偏角是磁北与地理北之间的角度差，而磁倾角是磁场线与地表的角度。
		
		\item \textbf{时间变化：} 地磁场会随时间发生缓慢的变化，包括偶极磁场的漂移和偶尔的磁极颠倒。
		
		\item \textbf{空间变化：} 地磁场还受到太阳风和地球磁层的影响，导致磁场在空间上发生变化。
		
		\item \textbf{保护作用：} 地磁场对地球生物和技术设施有重要的保护作用，特别是在防护来自太阳的宇宙射线方面。
	\end{enumerate}
	
	\begin{figure}[htbp]
		\centering
		\subfloat[]{
			\includegraphics[width=0.33\textwidth]{BA3GraAD3.jpg}
		}
		\subfloat[]{
			\includegraphics[width=0.33\textwidth]{BA3GraAD4.jpg}
		}
		\subfloat[]{
			\includegraphics[width=0.33\textwidth]{BA3GraAD5.jpg}
		}
		\caption{地磁场：强度、倾角、偏角（参考文献[1]）}
		\label{fig:figAD3}			
	\end{figure}
	
	这些特性使得地磁场成为航海导航、动物迁徙、地球科学研究和许多其他应用的关键因素。
	
	
	\subsection{实验前思考题}
	\begin{question}
		通过网上或图书馆查阅文献，列举某个AMR传感器在有关领域的应用实例，简要介绍其测量原理和方法?（2分）
	\end{question}
	AMR传感器是一种利用各向异性磁阻效应来测量磁场的角度或强度的磁敏元件。AMR传感器有许多应用实例，例如：
	\begin{enumerate}
		\item 开关检测：AMR传感器可以用于检测各种开关动作，如笔记本电脑的开关、冰箱的门开关、门锁的防盗装置等。AMR传感器的测量原理是，当外加磁场与电流方向夹角为0°或180°时，电阻最大；当夹角为90°时，电阻最小。通过检测电阻的变化，可以判断开关的状态。
		\item  位置检测：AMR传感器可以用于检测位置，如气缸的开关、电子锁的位置、液面的高度等。AMR传感器的测量方法是，将磁铁固定在被测物体上，将AMR传感器固定在参考位置上，当被测物体移动时，磁铁对AMR传感器的磁场影响会改变，从而改变AMR传感器的电阻，通过测量电阻的变化，可以计算出被测物体的位置。
		\item 旋转检测：AMR传感器可以用于检测旋转，如水表、速度表、风速计、轨迹球等。
		AMR传感器的测量方法是，将磁铁固定在旋转轴上，将AMR传感器固定在静止部分上，当旋转轴转动时，磁铁对AMR传感器的磁场角度会改变，从而改变AMR传感器的电阻，通过测量电阻的周期性变化，可以计算出旋转轴的转速和转向。
		\item 全面检测：AMR传感器可以用于检测大范围的磁场，如替换磁簧开关、检测各种设备的外部磁场等。AMR传感器的测量方法是，将AMR传感器放置在需要检测的区域内，当有磁场存在时，AMR传感器的电阻会发生变化，通过测量电阻的变化，可以判断磁场的存在和强度。
	\end{enumerate}
	
	\begin{figure}[htbp]
		\centering
		\includegraphics[width=0.3\textwidth]{BA3GraAD6.jpg}
		\caption{AMR传感器}
		\label{fig:figAD6}
	\end{figure}
	
	具体应用实例：各向异性磁阻（AMR）传感器，例如由Analog Devices开发的ADA4570，是高精度位置测量领域的一个重要应用实例。AMR传感器利用铁磁性材料的一个特性，即材料的电阻值依赖于磁化方向，这一现象最早由威廉·汤姆森（Lord Kelvin）在1851年左右发现。
	这些传感器特别适用于测量移动和旋转元件的位置。例如，在电机控制应用中，AMR传感器可用于无刷直流电机的转子位置测量和控制。它们通常用于旋转元件的偏轴或轴端配置中，其中传感器的正弦/余弦输出为每个磁极提供绝对信息。这种安排允许精确控制和测量转子的位置。
	AMR传感器的测量原理涉及使用两个惠斯通桥，其中一个相对于另一个旋转了45度。通过正弦和余弦函数计算角度，代表相对于传感器从0到180度的方向。这种方法允许在180度机械旋转中进行精确测量，对于偶极磁体，电气周期在360度旋转中重复两次。AMR传感器在饱和状态下工作，这意味着绝对场强不那么关键，当使用强磁体时，可以实现稳健的系统。
	
	\begin{question}
		常用的磁阻元件有哪几种? （2分）
	\end{question}
	常用的磁阻元件主要有以下几种类型：
	\begin{enumerate}
		\item 普通磁阻（MR）元件：这是最基本的类型，依赖于磁场对材料电阻值的影响。
		
		\item 各向异性磁阻（AMR）元件：这些元件的电阻值不仅取决于磁场的强度，还取决于磁场的方向。
		
		\item 巨磁阻（GMR）元件：这些元件利用不同材料层之间的磁性耦合，提供比MR和AMR元件更显著的电阻变化。
		
		\item 隧道磁阻（TMR）元件：它们使用了量子力学中的隧道效应，通常提供比GMR更高的电阻变化率。
		
		这些磁阻元件在不同的应用领域中具有重要作用，例如在传感器、数据存储和读取设备等方面。
	\end{enumerate}
	
	半导体磁阻元件：利用半导体材料的磁阻效应，即半导体材料的电阻随着外加磁场的变化而变化。半导体磁阻元件有高灵敏度、高分辨率、低温漂移等优点，但也有易受温度、电流、应力等因素影响的缺点。半导体磁阻元件的代表有磁敏电阻（Magnetoresistive, MR）、巨磁阻（Giant Magnetoresistive, GMR）、隧道磁阻（Tunneling Magnetoresistive, TMR）等。
	
	三极管词组元件：利用三极管的磁控特性，即三极管的电流放大倍数随着外加磁场的变化而变化。三极管词组元件有低功耗、高稳定性、高可靠性等优点，但也有低灵敏度、低分辨率、易受温度影响的缺点。三极管词组元件的代表有磁控三极管（Magnetic Transistor, MT）、磁控晶体管（Magnetic Field Effect Transistor, MAGFET）等。
	
	%实验记录	
	\clearpage
	\begin{table}
		\renewcommand\arraystretch{1.7}
		\centering
		\begin{tabularx}{\textwidth}{|X|X|X|X|}
			\hline
			专业： & 物理学 & 年级： & 2022级 \\
			\hline
			姓名： & 杨舒云 & 学号： & 22344020\\
			\hline
			室温： & 22摄氏度 & 实验地点： & 实验室511 \\
			\hline
			学生签名：& 杨舒云 & 评分： &\\
			\hline
			实验时间：& 2023/12/14 & 教师签名：&\\
			\hline
		\end{tabularx}
	\end{table}
	
	\section{BA3 应用AMR测量磁场分布 \quad\heiti 实验记录}
	\subsection{实验内容、步骤与结果}
	
	\subsubsection{磁阻传感器特性测量（10分）}
	\begin{enumerate}
		\item 简述本次实验方法及过程
		
		实验所用的磁阻传感器是AMR（各向异性磁电阻）传感器，它是由四个磁敏感的电阻组成的电桥，当外加磁场改变时，电阻的阻值也会随之变化，从而导致电桥输出电压的变化。
		
		实验中使用赫姆霍兹线圈作为外加磁场的源，它可以产生均匀的磁场，其磁感应强度与线圈的电流成正比。
		
		实验分为两个部分：
		\begin{enumerate}
			\item 第一部分是测量磁阻传感器的磁电转换特性，即输出电压与磁感应强度的关系。实验步骤如下：
			\begin{itemize}
				\item 将磁阻传感器放置在赫姆霍兹线圈的中心，使其磁敏感方向与磁场方向平行。
				\item 按表9-2（参考文献[2]）的数据，从300mA开始，逐步调小赫姆霍兹线圈的电流，记录相应的输出电压值。
				\item 切换电流换向开关，使赫姆霍兹线圈的电流反向，磁场和输出电压也将反向。
				\item 逐步调大反向电流，记录反向输出电压值。
				\item 每次切换电流方向后，必须按复位按键消磁，以消除磁阻传感器的磁滞现象。
				\item 根据输出电压和磁感应强度的数据，绘制磁电转换特性曲线，计算磁阻传感器的灵敏度，即曲线的斜率。
			\end{itemize}
			\item 第二部分是测量磁阻传感器的各向异性特性，即输出电压与磁场方向的关系。实验步骤如下：
			\begin{itemize}
				\item 将赫姆霍兹线圈的电流调节至200mA，保持磁场强度不变。
				\item 测量磁场方向与磁敏感方向一致时的输出电压，即最大输出电压。
				\item 松开线圈水平旋转锁紧螺钉，每次将赫姆霍兹线圈与传感器盒整体转动10度后锁紧，改变磁场方向。
				\item 松开传感器水平旋转锁紧螺钉，将传感器盒向相反方向转动10度，保持磁敏感方向不变。
				\item 记录输出电压数据于表9-3（参考文献[2]）中，直到转动一周为止。
				\item 根据输出电压和磁场方向的数据，绘制各向异性特性曲线，观察输出电压与磁场方向的夹角的余弦函数关系。
			\end{itemize}
		\end{enumerate}
		
		\item 实验数据记录（\cref{tab:tab1}与\cref{tab:tab2}）
		
		\begin{table}[h]
			\centering
			\caption{AMR磁电转换特性的测量}
			\label{tab:tab1}
			\resizebox{\textwidth}{!}{%
				\begin{tabular}{|c|c|c|c|c|c|c|c|c|c|c|c|c|c|}
					\hline
					线圈电流/mA & 300 & 250 & 200 & 150 & 100 & 50 & 0 & -50 & -100 & -150 & -200 & -250 & -300 \\
					\hline
					磁感应强度/Gauss & 6 & 5 & 4 & 3 & 2 & 1 & 0 & -1 & -2 & -3 & -4 & -5 & -6 \\
					\hline
					输出电压/V & 1.501 & 1.270 & 1.027 & 0.775 & 0.518 & 0.258 & 0.000 & -0.264 & -0.525 & -0.777 & -1.025 & -1.260 & -1.487 \\
					\hline
				\end{tabular}
			}%
		\end{table}
		
		\begin{table}[h]
			\centering
			\caption{AMR方向特性的测量（$B_0=4Gauss$）}
			\label{tab:tab2}
			\begin{tabular}{|c|c|}
				\hline
				夹角$\alpha$/degree & 输出电压/V \\
				\hline
				0 & 1.027 \\
				10 & 1.013 \\
				20 & 0.970 \\
				30 & 0.899 \\
				40 & 0.799 \\
				50 & 0.675 \\
				60 & 0.524 \\
				70 & 0.364 \\
				80 & 0.177 \\
				90 & -0.006 \\
				\hline
			\end{tabular}
		\end{table}
		
		
		\item 实验现象描述
		
		\textbf{实验现象一：测量磁阻传感器的磁电转换特性}
		\begin{itemize}
			\item 当赫姆霍兹线圈电流从 300mA 逐步减小到 0mA 时，磁阻传感器的输出电压也随之减小，呈现出一个下降的曲线。
			\item 当赫姆霍兹线圈电流从 0mA 逐步增大到 -300mA 时，磁阻传感器的输出电压也随之增大，呈现出一个上升的曲线。
			\item 当赫姆霍兹线圈电流为 0mA 时，磁阻传感器的输出电压为零，此时磁场为零，磁阻效应消失。
			\item 当赫姆霍兹线圈电流的绝对值较小时，磁阻传感器的输出电压与磁感应强度的平方成正比，呈现出一个二次函数的形状。
			\item 当赫姆霍兹线圈电流的绝对值较大时，磁阻传感器的输出电压与磁感应强度成正比，呈现出一个线性函数的形状。
		\end{itemize}
		
		\textbf{实验现象二：测量磁阻传感器的各向异性特性}
		\begin{itemize}
			\item 当赫姆霍兹线圈电流为 200mA 时，磁阻传感器的输出电压达到最大值，此时所测磁场方向与磁敏感方向一致，磁阻效应最明显。
			\item 当赫姆霍兹线圈与传感器盒整体旋转时，磁阻传感器的输出电压随之变化，呈现出一个周期性的波动，每旋转 180 度，输出电压的极值重复出现。
			\item 当赫姆霍兹线圈与传感器盒整体旋转 90 度时，磁阻传感器的输出电压为零，此时所测磁场方向与磁敏感方向垂直，磁阻效应消失。
			\item 当赫姆霍兹线圈与传感器盒整体旋转的角度为 $a$ 时，磁阻传感器的输出电压与所测磁场在磁敏感方向的投影成正比，即与磁感应强度的余弦值成正比。
		\end{itemize}
		
	\end{enumerate}
	
	\subsubsection{亥姆霍兹线圈磁场分布测量（10分）}
	\begin{enumerate}
		\item 简述本次实验方法及过程
		
		实验所用的赫姆霍兹线圈是由两个共轴的圆形线圈组成，线圈之间的距离等于线圈的半径，线圈上通以同向的交流电流，产生近似均匀的磁场。
		
		实验中使用霍尔元件作为磁场的探测器，它可以感应出与磁感应强度成正比的霍尔电压，通过测量霍尔电压的大小和方向，可以推算出磁场的分布。
		
		\textbf{实验分为两个部分：}
		
		\textbf{第一部分是测量赫姆霍兹线圈轴线上的磁场分布，即磁感应强度与轴线距离的关系。实验步骤如下：}
		\begin{itemize}
			\item 将霍尔元件放置在赫姆霍兹线圈的中心，使其磁敏感方向与磁场方向平行。
			\item 调节线圈的电流和频率，使磁场达到稳定的状态。
			\item 测量霍尔元件的输出电压，记录为中心点的磁感应强度。
			\item 沿着轴线方向，每隔一定的距离，移动霍尔元件，重复测量输出电压，记录为不同位置的磁感应强度。
			\item 根据输出电压和磁感应强度的数据，绘制磁场分布曲线，观察磁场的均匀性和变化规律。
		\end{itemize}
		
		\textbf{第二部分是测量赫姆霍兹线圈空间内的磁场分布，即磁感应强度与空间位置的关系。实验步骤如下：}
		\begin{itemize}
			\item 将霍尔元件放置在赫姆霍兹线圈的中心，使其磁敏感方向与磁场方向平行。
			\item 调节线圈的电流和频率，使磁场达到稳定的状态。
			\item 测量霍尔元件的输出电压，记录为中心点的磁感应强度。
			\item 沿着径向方向，每隔一定的距离，移动霍尔元件，重复测量输出电压，记录为不同位置的磁感应强度。
			\item 根据输出电压和磁感应强度的数据，绘制磁场分布曲线，观察磁场的均匀性和变化规律。
		\end{itemize}
				
		\item 实验数据记录（\cref{tab:tab3}与\cref{tab:tab4}）
		
		说明：理论分析表明，可认为在赫姆霍兹线圈中部较大的区域内，磁场方向沿轴线方向，磁场大小基本不变。按表9-5（参考文献[2]）数据改变磁阻传感器的空间位置，记录X方向的磁 场产生的电压 VX，测量赫姆霍兹线圈空间磁场分布。
		
		\begin{table}[h]
			\centering
			\caption{亥姆霍兹线圈轴向磁场（x向）分布测量}
			\label{tab:tab3}
			\begin{tabular}{|c|c|c|c|c|c|c|c|c|c|c|c|}
				\hline
				位置x & -0.5R & -0.4R & -0.3R & -0.2R & -0.1R & 0 & 0.1R & 0.2R & 0.3R & 0.4R & 0.5R \\
				\hline
				计算值B(x)/B0 & 0.946 & 0.975 & 0.992 & 0.998 & 1.000 & 1.000 & 1.000 & 0.998 & 0.992 & 0.975 & 0.946 \\
				测量值B(x)/V & 1.420 & 1.463 & 1.485 & 1.496 & 1.497 & 1.500 & 1.498 & 1.497 & 1.490 & 1.468 & 1.432 \\
				测量值B(x)/Gauss & 5.680 & 5.852 & 5.940 & 5.984 & 5.988 & 6.000 & 5.992 & 5.988 & 5.960 & 5.872 & 5.728 \\
				测量值B(x)/B0 & 0.947 & 0.975 & 0.990 & 0.997 & 0.998 & 1.000 & 0.999 & 0.998 & 0.993 & 0.979 & 0.955 \\
				\hline
			\end{tabular}
		\end{table}
		
		\begin{table}[h]
			\centering
			\caption{亥姆霍兹线圈空间磁场（x,y）分布测量}
			\label{tab:tab4}
			\begin{tabular}{|c|ccccccc|}
				\hline
				\diagbox{Y}{$U_x/V$}{X} & 0 & 0.05R & 0.10R & 0.15R & 0.20R & 0.25R & 0.30R \\
				\hline
				0 & 1.027 & 1.027 & 1.029 & 1.027 & 1.028 & 1.027 & 1.024 \\
				0.05R & 1.030 & 1.020 & 1.031 & 1.030 & 1.030 & 1.029 & 1.025 \\
				0.10R & 1.028 & 1.028 & 1.029 & 1.029 & 1.020 & 1.029 & 1.025 \\
				0.15R & 1.027 & 1.027 & 1.028 & 1.028 & 1.030 & 1.029 & 1.028 \\
				0.20R & 1.026 & 1.028 & 1.029 & 1.031 & 1.032 & 1.033 & 1.033 \\
				0.25R & 1.027 & 1.028 & 1.028 & 1.030 & 1.033 & 1.035 & 1.037 \\
				0.30R & 1.025 & 1.025 & 1.027 & 1.029 & 1.034 & 1.041 & 1.045 \\
				\hline
			\end{tabular}
		\end{table}
				
		\item 实验现象描述
		
		\textbf{实验现象一：赫姆霍兹线圈轴线上的磁场分布测量}
		\begin{itemize}
			\item 当传感器磁敏感方向与赫姆霍兹线圈轴线一致，位置调节至赫姆霍兹线圈中心 ($X=0$) 时，输出电压达到最大值，此时磁感应强度为 $B_0$，由公式 (9-5) 可计算得到。
			\item 当传感器沿轴线平移时，输出电压随之变化，呈现出一个对称的波峰波谷的曲线，每平移 $0.1R$，输出电压的极值重复出现。
			\item 当传感器平移的距离 $X$ 越大时，输出电压的幅度越小，磁感应强度 $B(X)$ 也越小，由公式 (9-4)（参考文献[2]） 可计算得到。
			\item 当传感器平移的距离 $X$ 等于线圈半径 $R$ 时，输出电压为零，磁感应强度 $B(X)$ 也为零，此时两个线圈在该点产生的磁场相互抵消。
		\end{itemize}
		
		\textbf{实验现象二：赫姆霍兹线圈空间磁场分布测量}
		\begin{itemize}
			\item 当传感器磁敏感方向与赫姆霍兹线圈轴线一致，位置调节至赫姆霍兹线圈中心 ($X=0$, $Y=0$) 时，输出电压达到最大值，此时磁感应强度为 $B_0$，由公式 (9-5) 可计算得到。
			\item 当传感器沿 $X$ 方向或 $Y$ 方向移动时，输出电压随之变化，呈现出一个类似于高斯函数的曲线，移动的距离越大，输出电压的幅度越小，磁感应强度 $B(X)$ 或 $B(Y)$ 也越小。
			\item 当传感器移动的距离 $X$ 或 $Y$ 小于 $0.2R$ 时，输出电压的变化幅度很小，磁感应强度 $B(X)$ 或 $B(Y)$ 与 $B_0$ 的相对误差小于百分之一，说明赫姆霍兹线圈在中心附近的区域内，磁场较为均匀。
			\item 当传感器移动的距离 $X$ 或 $Y$ 大于 $0.2R$ 时，输出电压的变化幅度较大，磁感应强度 $B(X)$ 或 $B(Y)$ 与 $B_0$ 的相对误差超过百分之一，说明赫姆霍兹线圈在边缘的区域内，磁场较为不均匀。
		\end{itemize}
		
	\end{enumerate}
	
	\subsubsection{地磁场测量（10分）}
	\begin{enumerate}
		\item 简述本次实验方法及过程
		
		实验所用的磁阻传感器是一种能够感应外加磁场的电阻，当磁场变化时，电阻的阻值也会随之变化，从而导致电桥输出电压的变化。
		
		实验中使用赫姆霍兹线圈作为外加磁场的源，它可以产生均匀的磁场，其磁感应强度与线圈的电流成正比。
		
		实验分为以下几个步骤：
		\begin{itemize}
			\item 将磁阻传感器放置在赫姆霍兹线圈的中心，使其磁敏感方向与线圈轴线垂直。
			\item 将赫姆霍兹线圈的电流和补偿电流调节至零，使外加磁场为零。
			\item 调节传感器盒上平面与仪器底板平行，用水准气泡盒检查传感器是否水平。
			\item 松开线圈水平旋转锁紧螺钉，在水平面内仔细调节传感器方位，使输出电压最大。此时，传感器磁敏感方向与地理南北方向的夹角即为磁偏角，记录此角度。
			\item 松开传感器绕轴旋转锁紧螺钉，在垂直面内调节磁敏感方向，使输出电压最大。此时，传感器磁敏感方向与水平面的夹角即为磁倾角，记录此角度。
			\item 在自己创作的数据表格中，记录输出最大时的输出电压值 $U1$ 后，将传感器水平旋转 $180$ 度，记录此时的输出电压 $U2$，计算 $U = (U1 - U2)/2$ 作为地磁场磁感应强度的测量值。
			\item 重复以上步骤，改变传感器方位，测量不同方向的磁场分量，计算地磁场的大小和方向。
		\end{itemize}
		
		
		\item 实验数据记录（\cref{fig:fig1}）
		
		\begin{figure}[htbp]
			\centering
			\includegraphics[width=0.5\textwidth]{BA3Gra1.png}
			\caption{地磁场测量}
			\label{fig:fig1}
		\end{figure}
		
		在实验室内测量地磁场时，建筑物的钢筋分布，同学携带的铁磁物质，都可能影响测量 结果，因此，此实验重在掌握测量方法。
		
		\item 实验现象描述
		
		如\cref{tab:tab5}所示。
		
		\begin{table}[h]
			\centering
			\caption{地磁场测量}
			\label{tab:tab5}
			\begin{tabular}{|c|c|}
				\hline
				倾角/degree & $U_1/V$ \\
				\hline
				水平 & 0.152 \\
				40 & 0.164 \\
				30 & 0.165 \\
				20 & 0.164 \\
				\hline
			\end{tabular}
		\end{table}
		
		用磁阻传感器测量地磁场的大小和方向，需要先用亥姆霍兹线圈产生一个已知的磁场，校准传感器的灵敏度。
		传感器的输出电压与磁场的分量成正比，当传感器的磁敏感方向与磁场的分量平行时，输出电压最大，当垂直时，输出电压为零。
		调节传感器的方位，使其与地磁场的水平分量平行，记录输出电压，再旋转180度，记录输出电压，计算地磁场的水平分量和磁偏角。
		调节传感器的倾角，使其与地磁场的垂直分量平行，记录输出电压，计算地磁场的垂直分量和磁倾角。
		用输出电压的差值除以二，消除电桥偏差的影响，计算地磁场的总强度。
		避免周围的铁磁物质对测量结果的干扰。
		
		除此之外，我还做了以下观察：我将正在接通电话的手机靠近AMR传感器，在地磁场的基础上，磁感应强度有显著变化。
		
	\end{enumerate}
	
	
	\subsection{原始数据记录}
	实验记录本上的原始数据见\cref{fig:figA1}。
	预习报告签字见\cref{fig:figA3}。
	实验台桌面整理见\textbf{附件}部分(\cref{fig:figA4})。
	
	\begin{figure}[htbp]
		\centering
		\subfloat[]{
			\includegraphics[width=0.5\textwidth]{BA3GraA1.jpg}
		}
		\subfloat[]{
			\includegraphics[width=0.5\textwidth]{BA3GraA2.jpg}
		}
		\caption{原始数据记录}
		\label{fig:figA1}			
	\end{figure}
	
	
	
%	\subsection{实验过程中遇到的问题记录}
%	\begin{enumerate}
%		\item 
%	\end{enumerate}
	
	
	
	%分析与讨论	
	\clearpage
	\begin{table}
		\renewcommand\arraystretch{1.7}
		\begin{tabularx}{\textwidth}{|X|X|X|X|}
			\hline
			专业：& 物理学 &年级：& 2022级\\
			\hline
			姓名： & 杨舒云 & 学号：& 22344020\\
			\hline
			日期：& 2023/12/14 & 评分： &\\
			\hline
		\end{tabularx}
	\end{table}
	
	\section{BA3 应用AMR测量磁场分布 \quad\heiti 分析与讨论}
	\subsection{实验数据分析}
	\subsubsection{磁阻传感器特性测量（12分）}
	\begin{enumerate}
		\item 实验数据处理，根据原始测量数据，以磁感应强度为横轴，输出电压为纵轴，将\cref{fig:fig1}数据作图。
		
		根据要求，我用Python语言实现了作图，代码详见附录。如\cref{fig:fig2}所示。
		
		\begin{figure}[htbp]
			\centering
			\includegraphics[width=0.7\textwidth]{BA3Gra2.png}
			\caption{AMR磁电转换特性曲线}
			\label{fig:fig2}
		\end{figure}
		
		\item 分析上图，确定所用传感器的线性工作范围及灵敏度。
		
		事实上，对于本实验的数据处理来说，利用Origin远比Python方便，因此，我有用Origin进行了一遍重复工作，如\cref{fig:fig3}所示。
		
		\begin{figure}[htbp]
			\centering
			\includegraphics[width=0.7\textwidth]{BA3Gra3.jpg}
			\caption{磁电转换特性曲线（Origin）}
			\label{fig:fig3}
		\end{figure}
		
		根据图像，分析如下：
		\begin{itemize}
			\item 所用传感器的线性工作范围大约是从 $-6$ Gauss 到 $6$ Gauss，因为在这个范围内，输出电压与磁感应强度呈现很好的线性关系。
			\item 所用传感器的灵敏度可以用拟合曲线的斜率表示，即 $a = \frac{\Delta V}{\Delta B}$，其中 $\Delta V$ 是输出电压的变化量，$\Delta B$ 是磁感应强度的变化量。根据拟合结果，$a \approx 0.248$，即传感器的灵敏度约为 $0.25302 \pm 0.00121$ V/Gauss。
			\item 事实上，根据\textbf{实验现象}所述，当赫姆霍兹线圈电流的绝对值较小时，磁阻传感器的输出电压与磁感应强度的平方成正比，呈现出一个二次函数的形状。
			当赫姆霍兹线圈电流的绝对值较大时，磁阻传感器的输出电压与磁感应强度成正比，呈现出一个线性函数的形状。
		\end{itemize}		
		
		\item 根据原始测量数据，以夹角$\alpha$为横轴，输出电压为纵轴，将\cref{tab:tab2}数据作图，判断曲线有何规律。
		
		首先，我用Origin进行了初步的数据处理，并采用二次函数进行拟合，如\cref{fig:fig4}所示。可以看出，相关性已经非常好了，因此我们可以认为输出电压与夹角有着二次函数关系。
		
		\begin{figure}[htbp]
			\centering
			\includegraphics[width=0.7\textwidth]{BA3Gra4.jpg}
			\caption{AMR方向特性曲线（Origin）}
			\label{fig:fig4}
		\end{figure}
		
		进一步，我考虑从实验原理出发，利用Python语言进行了进一步地探究，如\cref{fig:fig5}所示。
		
		\begin{figure}[htbp]
			\centering
			\includegraphics[width=0.7\textwidth]{BA3Gra5.png}
			\caption{AMR方向特性曲线}
			\label{fig:fig5}
		\end{figure}
		
		从散点图可以看出，数据呈现出一个\textbf{下降的抛物线}的趋势，随着夹角 $\alpha$ 的增大，输出电压 $V$ 逐渐减小。从拟合曲线可以看出，拟合函数能够较好地拟合数据，拟合参数的置信区间也较小，说明拟合效果较好。拟合函数的形式为：
		\[ V = 0.512 \cdot \cos(\alpha)^2 - 0.001 \cdot \cos(\alpha) + 0.005 \]
		这个函数表明，输出电压 $V$ 与夹角 $\alpha$ 的余弦的二次函数有关，这可能与实验装置的原理有关。
		
		曲线是一个\textbf{下降的抛物线}，这说明输出电压与磁场方向的夹角的余弦的平方成正比，即 $V \propto \cos^2(\alpha)$，其中 $\alpha$ 是磁场方向与磁敏感方向的夹角。
		这个关系可以用磁阻传感器的原理来解释，它是由四个相同的磁阻元件构成的惠斯通电桥，每个元件的电阻会随着磁场方向的改变而改变，从而影响电桥的输出电压。
		当磁场方向与磁敏感方向一致时，即 $\alpha = 0$，电桥的输出电压达到最大值，因为此时磁阻元件的电阻差异最大，电桥的不平衡程度最高。
		当磁场方向与磁敏感方向垂直时，即 $\alpha = 90$，电桥的输出电压接近于零，因为此时磁阻元件的电阻差异最小，电桥的不平衡程度最低。
		当磁场方向在两者之间变化时，电桥的输出电压随着 $\cos^2(\alpha)$ 的变化而变化，呈现出抛物线的形状。		
		
	\end{enumerate}
	
	\subsubsection{亥姆霍兹线圈磁场分布测量（12分）}
	\begin{enumerate}
		\item 实验数据处理，根据\cref{tab:tab3}原始测量数据作图。
		
		根据要求，我用Python语言实现了作图，代码详见附录。如\cref{fig:fig6}所示（为了方便比对，我还作了理论值图像）。
		
		\begin{figure}[htbp]
			\centering
			\includegraphics[width=0.8\textwidth]{BA3Gra6.png}
			\caption{亥姆霍兹线圈轴向磁场（x向）分布图}
			\label{fig:fig6}
		\end{figure}
		
		\item 根据上图分析亥姆霍兹线圈的轴向磁场分布特点。
		
		在分析亥姆霍兹线圈的轴向磁场分布时，主要关注的是磁场强度随位置变化的规律。赫姆霍兹线圈设计的目的是在其中心区域产生一个尽可能均匀的磁场。
		
		从数据和作图来看，我们可以观察到以下特点：
		
		\begin{itemize}
			\item \textbf{均匀性：} 图中的“计算值”曲线（红色）在中心区域（从$-0.1R$到$0.1R$）非常平坦，说明理论上磁场在这个区间非常均匀。这与亥姆霍兹线圈的预期特性一致。实际的“测量值”曲线（蓝色）也显示了相似的趋势，但在中心区域以外的均匀性下降更明显。
			
			\item \textbf{磁场强度：} 磁场强度在中心最大，即$B(x)/B_0$为1，或者$B(x)$达到测量的最大值。这是因为亥姆霍兹线圈的设计就是为了在中心点产生最大的磁场强度。
			
			\item \textbf{对称性：} 在理论计算和实际测量中，磁场分布曲线关于中心点对称，这反映了线圈几何结构的对称性。
			
			\item \textbf{磁场衰减：} 随着距离中心点的增加，磁场强度逐渐减小。这一点在计算值和测量值中都有所体现，尽管测量值的减小趋势略微不同。
			
			\item \textbf{测量与计算的差异：} 理论计算的曲线与实际测量的曲线在形状上大致相似，但具体数值有所差异。这可能是因为实验中的环境干扰、设备精度、线圈制造的不完美等因素导致。
		\end{itemize}
		
		亥姆霍兹线圈的轴向磁场分布的理想特点是在中心区域非常均匀，适合于需要均匀磁场环境的实验。然而，实际应用中，由于多种因素的影响，实际测量值可能会与理论计算值有所偏差，因此在使用过程中需要对这些因素进行校准和考虑。
		
		
		\item 根据\cref{tab:tab4}数据，选择合理的方法（绘图或计算）讨论亥姆赫兹线圈的空间磁场分布特点。
		
		根据要求，我用Python语言实现了作图，代码详见附录。如\cref{fig:fig7}所示。进一步，我利用对称性，将分布图推广到了全空间的情况，如\cref{fig:fig8}所示。
		
		\begin{figure}[htbp]
			\centering
			\includegraphics[width=0.7\textwidth]{BA3Gra7.png}
			\caption{亥姆霍兹线圈空间磁场（x,y）分布图（局域）}
			\label{fig:fig7}
		\end{figure}
		
		\begin{figure}[htbp]
			\centering
			\subfloat[]{
				\includegraphics[width=0.5\textwidth]{BA3Gra8.png}
			}
			\subfloat[]{
				\includegraphics[width=0.5\textwidth]{BA3Gra9.png}
			}
			\caption{亥姆霍兹线圈空间磁场（x,y）分布}
			\label{fig:fig8}			
		\end{figure}
		
		从图表和数据中我们可以看到：
		
		\begin{itemize}
			\item \textbf{中心区域均匀性：} 数据表明，在中心区域（$X$ 和 $Y$ 都小于 $0.2R$ 的范围内），电压值变化不大，这说明磁场在这个区域较为均匀。这与理论分析是一致的，理论上在这个区域内$(B_X-B_0)/B_0$ 应小于百分之一，$B_Y/B_X$ 应小于万分之二，意味着磁场方向主要沿轴线方向，磁场大小基本不变。
			
			\item \textbf{磁场方向：} 实验步骤中提到，磁场的方向主要沿轴线方向。这与毕奥-萨伐尔定律计算的轴对称特性相符合。
			
			\item \textbf{电压变化趋势：} 当 $X$ 和 $Y$ 的值增大时，特别是在 $0.25R$ 到 $0.30R$ 的区间内，我们可以看到电压值有显著的增加，这可能意味着在这些较远的位置，磁场强度有所下降。
			
			\item \textbf{空间磁场非均匀性：} 随着距离中心点的增加，电压值的变化更加显著，特别是在图表的边缘。这表明磁场在远离中心的区域变得不那么均匀。
			
			\item \textbf{磁场衰减：} 在距离中心较远的地方（尤其是在 $Y$ 轴上），磁场强度有明显的衰减，这反映在电压值的增加上。
		\end{itemize}
		
		综上所述，赫姆霍兹线圈的设计确实在中心区域提供了一个相对均匀的磁场，但均匀性随着距离中心的增加而下降。这些观察结果对于需要在均匀磁场中进行精确测量的实验来说非常重要，例如量子物理实验和精密磁共振成像。在实际应用中，需要对这些非均匀性进行校准，以确保实验结果的准确性。		
		
	\end{enumerate}
	
	\subsubsection{地磁场测量（4分）}
	通过网上或图书馆查阅文献，简要介绍地磁场可能的产生机理?
	
	地磁场的产生机理主要基于地球内部的动力学过程，这一过程被称为地球的“地幔对流”。以下是目前科学界普遍认同的几种主要机理：
	
	\begin{itemize}
		\item 地球动力发电机理论：这是最广泛接受的解释，认为地磁场是由地球内核的流体动力学运动产生的。地球的外核主要由铁和镍组成，处于液态。这些金属在地球自转过程中流动，产生电流，进而产生磁场。
		
		\item \textbf{热对流效应}：地球内部的热对流也被认为是地磁场产生的一个重要因素。地球内部的热能导致金属液体发生对流，这种对流运动带动金属的电荷移动，产生电流，进而产生磁场。
		
		\item \textbf{旋转和对流相结合的效应}：地球的旋转与外核液体的对流相结合，可能导致复杂的流体动力学效应，这些效应在宏观上表现为磁场的产生。
		
		\item \textbf{化学反应}：地核的某些化学反应可能也对磁场的产生有一定的贡献，尽管这一点还没有得到广泛的科学共识。
		
		\item \textbf{自发磁化效应}：有研究指出，地球内部某些物质可能具有自发磁化的特性，这些物质在地球形成初期可能对地磁场的产生有所贡献。
	\end{itemize}
	
	这些理论都指出，地磁场是一个动态变化的系统，其变化与地球内部的物理条件和化学成分密切相关。对于这些机理的进一步探究，需要依靠地球物理学、地球化学和流体动力学等多个学科的深入研究。同时，利用卫星和地面观测数据，可以帮助科学家更好地理解这些复杂过程。
	
	除以上外，我还调研了相关研究的最新进展：
	\begin{itemize}
		\item \textbf{地核发电机的可持续性}：由外核流体运动产生的地球磁场已运行超过34亿年。然而，在如此长的时间内维持地核发电机的确切机制仍存在争议。提出的机制包括岁差、潮汐和对流。由热化学对流驱动的地核发电机的数值模型可以解释现有地核发电机的大多数观测特性。对于维持地核发电机，对流的作用在很大程度上取决于地核的热导率，这仍然是一个争议的话题。如果地核的热导率低于100 W m⁻¹ K⁻¹，对流可以为地核发电机提供足够的动力。地核中轻元素的析出可能进一步增加对流维持地核发电机的热导率上限。
		
		\item \textbf{早期磁场的强度和起源}：罗切斯特大学的研究表明，地球早期的磁场比先前认为的要强。这一发现基于锆石晶体的新数据，表明大约在40亿年前，地球就有了一个强磁场，其动力源于一种不同的机制，因为内核还没有形成。一个提出的机制是地球内部镁氧化物的化学沉淀，与形成地球月球的巨大撞击有关。这一过程可能在内核形成之前推动了对流和地核发电机的运行。在镁氧化物耗尽后，磁场几乎在5.65亿年前崩溃，但内核的形成为地核发电机提供了一个新的能源，维持了地球当前的磁场屏障。
		
		\item \textbf{关于地磁场最早存在的争论}：在澳大利亚西部一处古老的露头中，对微观矿物的研究，特别是锆石，最初暗示了地球磁场的痕迹可能可以追溯到42亿年前。然而，最近麻省理工学院领导的研究发现这些锆石不可靠，无法记录古老磁场。这项研究挑战了地球磁场存在于35亿年前之前的观点。磁场形成的时间对于理解产生机制至关重要，特别是因为地球核心据信只在10亿年前开始凝固，这表明在此之前磁场是由其他机制驱动的。
		
		\item \textbf{对早期生命的影响}：地球磁场的起源和早期存在对于地球上生命首次出现的环境具有重要意义。理解磁场形成的时间和机制，可以阐明支持地球前十亿年生命出现的条件。
	\end{itemize}
	
	
	\clearpage
	\subsection{实验后思考题}
	\begin{question}
		调研精确测量磁场的其他方法，与磁阻效应相比，有何优劣? 
	\end{question}
	\textbf{精确测量磁场的方法除了磁阻效应（如磁阻抗传感器），还包括多种其他技术。这些技术包括：}
	\begin{itemize}
		\item \textbf{SQUID (超导量子干涉装置)}：利用超导环的量子性质测量非常微弱的磁场。
		\item \textbf{感应式拾取线圈}：通过感应原理测量磁场变化。
		\item \textbf{VSM (振动样品磁强计)}：测量样品因磁场而产生的振动。
		\item \textbf{脉冲场提取磁强计}：适用于高强度磁场的测量。
		\item \textbf{扭矩磁强计}：测量磁场对物体产生的扭矩。
		\item \textbf{法拉第力磁强计}：利用磁场对物体产生的力进行测量。
		\item \textbf{光学磁强计}：使用光的性质来测量磁场。
	\end{itemize}
	
	\textbf{与磁阻效应相比，这些方法各有优缺点：}
	\begin{itemize}
		\item \textbf{灵敏度和精度}：SQUID是目前最灵敏的磁场测量设备，适用于极低强度的磁场测量。磁阻传感器通常灵敏度较低，但足以满足多数工业应用需求。
		\item \textbf{温度依赖性}：SQUID对温度非常敏感，需要低温环境。磁阻传感器对温度的依赖性较小。
		\item \textbf{成本和复杂性}：SQUID设备成本高，运行复杂。磁阻传感器相对简单、成本较低。
		\item \textbf{应用范围}：不同类型的磁强计适用于不同的应用领域。例如，SQUID在医学成像和地球物理学中非常有用，而磁阻传感器广泛用于工业和消费电子产品。
	\end{itemize}
	
	\begin{question}
		对提高测量精度，你有何改进的建议?
	\end{question}
	要提高使用各向异性磁阻(AMR)传感器测量磁场分布的精度，可以考虑以下几个方面的改进措施：
	\begin{itemize}
		\item 增强环境控制：确保实验环境中的外部磁场干扰最小。例如，远离电器设备和金属物体，使用屏蔽材料减少外部磁场的影响。
		
		\item \textbf{校准和调整传感器}：定期对AMR传感器进行校准，确保其精确性。此外，调整传感器的方位和位置，使其与被测磁场的方向更加一致，从而提高测量精度。
		
		\item \textbf{优化电桥电路}：仔细检查和优化电桥电路，确保所有电阻值精确匹配，以减少系统误差。
		
		\item \textbf{使用高质量的电源}：使用稳定、无干扰的电源为实验设备供电，减少电源波动对测量结果的影响。
		
		\item \textbf{增强数据处理能力}：采用高级的数据处理技术，如数字滤波和信号平滑，来提高测量数据的质量和可靠性。
		
		\item \textbf{采用温度补偿技术}：考虑温度对AMR传感器性能的影响，并采用相应的温度补偿措施。
		
		\item \textbf{提高实验技术水平}：确保实验操作人员对实验设备和方法有充分的了解和熟练的操作技能。
		
		\item \textbf{复位/反向置位功能的有效利用}：在遇到磁畴排列紊乱时，正确使用复位/反向置位功能，恢复传感器的特性。
	\end{itemize}
	
	这些措施可以显著提高AMR传感器在测量磁场分布时的准确性和可靠性。
	
	
	\clearpage
	\section{BA3 应用AMR测量磁场分布 \quad\heiti 参考文献与附件}
	\subsection{参考文献}
	[1] 维基百科.
	
	[2] 实验讲义.
	
	[3] Earth’s early magnetic field was stronger than previously believed.
	
	[4] The origin and evolution of Earth’s ancient magnetic field.
	
	[5] The early geodynamo: How old is Earth’s magnetic field?
	
	\subsection{附件}
	预习报告签字：
	
	\begin{figure}[htbp]
		\centering
		\includegraphics[width=0.5\textwidth]{BA3GraA3.jpg}
		\caption{预习报告签字}
		\label{fig:figA3}
	\end{figure}
	
	实验台桌面整理：
	
	\begin{figure}[htbp]
		\centering
		\subfloat[]{
			\includegraphics[width=0.4\textwidth]{BA3GraA4.jpg}
		}
		\subfloat[]{
			\includegraphics[width=0.4\textwidth]{BA3GraA5.jpg}
		}
		\caption{试验台桌面整理}
		\label{fig:figA4}			
	\end{figure}
	
	\clearpage
	代码记录：
	\begin{lstlisting}[style=pythonstyle,caption=代码记录1]
		# 导入必要的库
		import numpy as np
		import matplotlib.pyplot as plt
		from scipy.optimize import curve_fit
		import pandas as pd
		from mpl_toolkits.mplot3d import Axes3D
		
		## AMR磁电转换特性
		# 定义数据
		x = np.array([6, 5, 4, 3, 2, 1, 0, -1, -2, -3, -4, -5, -6]) # 磁感应强度
		y = np.array([1.501, 1.270, 1.027, 0.775, 0.518, 0.258, 0.000, -0.264, -0.525, -0.777, -1.025, -1.260, -1.487]) # 输出电压
		
		# 定义线性函数
		def linear(x, a, b):
		return a * x + b
		
		# 对数据进行拟合，得到最优参数a和b
		popt, pcov = curve_fit(linear, x, y)
		a = popt[0]
		b = popt[1]
		
		# 计算拟合曲线的y值
		y_fit = linear(x, a, b)
		
		# 绘制散点图
		plt.scatter(x, y, label='data', color='blue', marker='o')
		
		# 绘制折线图
		#plt.plot(x, y, label='line', color='red', linestyle='-')
		
		# 绘制样条连接图
		#plt.plot(x, y, label='spline', color='green', linestyle='--')
		
		# 绘制拟合曲线
		plt.plot(x, y_fit, label='fit', color='black', linestyle='-.')
		
		# 设置图例
		plt.legend()
		
		# 设置标题
		plt.title('Output voltage vs. magnetic induction intensity')
		
		# 设置坐标轴标签
		plt.xlabel('Magnetic induction intensity/Gauss')
		plt.ylabel('Output voltage/V')
		
		# 显示图像
		plt.show()
		
		## AMR各向异性
		# 将数据转换成DataFrame对象
		data = pd.DataFrame({"alpha": [0, 10, 20, 30, 40, 50, 60, 70, 80, 90], "V": [1.027, 1.013, 0.970, 0.899, 0.799, 0.675, 0.524, 0.364, 0.177, -0.006]})
		
		# 绘制散点图
		plt.scatter(data["alpha"], data["V"], label="data")
		plt.xlabel("alpha/degree")
		plt.ylabel("U/V")
		plt.title("AMR directional characteristics")
		
		# 定义拟合函数，关于夹角余弦的二次函数
		def fit_func(x, a, b, c):
		return a * np.cos(x * np.pi / 180) ** 2 + b * np.cos(x * np.pi / 180) + c
		
		# 对数据进行拟合，得到拟合参数
		popt, pcov = curve_fit(fit_func, data["alpha"], data["V"])
		a, b, c = popt
		print(f"The fitting function is: V = {a:.3f} * cos(alpha) ** 2 + {b:.3f} * cos(alpha) + {c:.3f}")
		
		# 绘制拟合曲线
		x = np.linspace(0, 90, 100)
		y = fit_func(x, a, b, c)
		plt.plot(x, y, label="fit", color='black')
		plt.legend()
		plt.show()
		
		## 亥姆霍兹轴向分布
		# 定义横坐标和纵坐标的数据
		x = [-0.5, -0.4, -0.3, -0.2, -0.1, 0, 0.1, 0.2, 0.3, 0.4, 0.5] # 位置x
		y1 = [0.946, 0.975, 0.992, 0.998, 1.000, 1.000, 1.000, 0.998, 0.992, 0.975, 0.946] # 计算值B(x)/B0
		y2 = [5.68, 5.852, 5.94, 5.984, 5.988, 6, 5.992, 5.988, 5.96, 5.872, 5.728] # 测量值B(x)/Gauss
		
		# 创建一个新的图形
		plt.figure()
		
		# 在第一幅子图中绘制计算值B(x)/B0与位置x的关系
		plt.subplot(2, 1, 1) # 两行一列的第一幅
		plt.plot(x, y1, 'r-') # 红色实线
		plt.xlabel('x') # 横坐标标签
		plt.ylabel('B(x)/B0') # 纵坐标标签
		plt.title('Calculated value') # 子图标题
		
		# 在第二幅子图中绘制测量值B(x)/Gauss与位置x的关系
		plt.subplot(2, 1, 2) # 两行一列的第二幅
		plt.plot(x, y2, 'b-') # 蓝色虚线
		plt.xlabel('x') # 横坐标标签
		plt.ylabel('B(x)/Gauss') # 纵坐标标签
		plt.title('Measured value') # 子图标题
		
		# 调整子图之间的间距
		plt.tight_layout()
		
		# 显示图形
		plt.show()
		
		## 亥姆霍兹空间分布
		# 创建数据
		x = np.array([0, 0.05, 0.1, 0.15, 0.2, 0.25, 0.3])
		y = np.array([0, 0.05, 0.1, 0.15, 0.2, 0.25, 0.3])
		z = np.array([[1.027, 1.027, 1.029, 1.027, 1.028, 1.027, 1.024],
		[1.030, 1.020, 1.031, 1.030, 1.030, 1.029, 1.025],
		[1.028, 1.028, 1.029, 1.029, 1.020, 1.029, 1.025],
		[1.027, 1.027, 1.028, 1.028, 1.030, 1.029, 1.028],
		[1.026, 1.028, 1.029, 1.031, 1.032, 1.033, 1.033],
		[1.027, 1.028, 1.028, 1.030, 1.033, 1.035, 1.037],
		[1.025, 1.025, 1.027, 1.029, 1.034, 1.041, 1.045]])
		
		# 创建网格
		X, Y = np.meshgrid(x, y)
		
		# 创建图形
		fig = plt.figure()
		ax = fig.add_subplot(111, projection='3d')
		
		# 绘制曲面
		#surf = ax.plot_surface(X, Y, z, cmap='coolwarm', edgecolor='black', linewidth=0.5)
		surf = ax.plot_surface(X, Y, z, cmap='plasma', edgecolor='white', linewidth=0.5)
		
		# 添加标题和图例
		#ax.set_title('三维分布图')
		#fig.colorbar(surf, shrink=0.5, aspect=5, label='$U_x/V$')
		
		# 设置坐标轴标签
		ax.set_xlabel('X/R')
		ax.set_ylabel('Y/R')
		ax.set_zlabel('$U_x/V$')
		
		# 调整视角和比例
		ax.view_init(30, -45)
		ax.set_zlim(1.02, 1.05)
		
		# 显示图形
		plt.show()
	\end{lstlisting}
	
	
\end{document}